\chapter{Translation: Transport or Rewriting?}
\section{Four Characters on a Cola Can}
First, pick up a can of cola.

You need not actually buy one; hold it in your imagination: a red can, a white ribbon, just the right size for one hand. Now turn it around to the side printed in Chinese. Four characters: \textbf{Kekou Kele}, the Chinese name for Coca-Cola.

Look at those four characters for ten seconds. They form the most widely circulated translation in China, and possibly the most successful translation in the world.

The original English name, ``Coca-Cola'', comes from ingredients: coca leaves and kola nuts, the names of two plants. If translation simply transports meaning from one language to another, this drink's Chinese name ought to be something like ``coca-leaf-and-kola-nut soda''. Nobody would buy something with that name.

Kekou Kele takes another route. First it preserves the sound: ke, kou, ke, le, four syllables whose outline roughly matches Coca-Cola. Then it secretly does something a carrier should not do: it slips in a layer of meaning absent from the original. Kekou means palatable; kele means joyful. Take a sip and feel refreshed and happy: that is advertising, not translation. Yet the name has been printed on the can for decades, entering the lives of generations. Various stories circulate about its origin, including one about a public naming contest; their details are difficult to verify individually. Another popular anecdote says an earlier translation was bizarre enough to frighten customers. Its truth is equally uncertain, so treat it as an anecdote. We reserve judgement on those stories, but the four characters on the can are certain: a translation can be so good that people forget it is a translation.

But wait. Turn the can again and read the ingredients: water, high-fructose corn syrup, white sugar, food additives\ldots This side is translated too, and every word must be right. ``Water'' cannot become ``joy''; ``sugar'' cannot become ``happiness''. Printing it incorrectly can bring legal liability.

On the same can, on the same sheet of aluminium, live two completely opposite kinds of translation: one rewrites the original meaning so successfully that it disappears; the other guards that meaning so strictly that not a word may go wrong.

Which is ``real translation''? Is translation moving goods from one box into another, or remaking the goods altogether? Starting from this can, this chapter will pursue that question. It is older and more dangerous than you might imagine, because the same choice arises not only on drink cans but in Buddhist scriptures, treaties, operating theatres, and the translation software on your phone.

\section{A Translation Hall in Chang'an}
Now come with me to another scene.

The time: a winter dawn in the middle of the seventh century CE. The place: the Great Ci'en Monastery in Chang'an, capital of the Tang dynasty. Day has not quite broken. The monastery bell has sounded, cold rises from the ground, and a group of monks crosses the courtyard towards a large room. Lamps and candles are already lit inside. Long desks stand in a row, bearing not sheets of paper but treated leaves: palm-leaf manuscripts brought from India and the Western Regions, covered in Sanskrit.

Seated at the centre is Xuanzang. Yes, the historical model for the monk you know from \emph{Journey to the West}. But here there is no monkey or pig, only a monk who has spent more than a decade crossing deserts and snowy mountains and brought over 650 Buddhist works back from India. Fetching the scriptures was only the first half. The second, longer half was translating them into Chinese. He worked at it for nearly twenty years, until his death, translating over seventy works in more than 1,300 fascicles.

Translating scripture was not a solitary person silently writing in a study. The Tang official translation hall was a production line with strictly divided responsibilities. Xuanzang himself was the ``translation master'', holding the Sanskrit text, reading it sentence by sentence, and explaining its meaning. A ``recorder'' wrote down his spoken Chinese explanation; a ``verifier of meaning'' checked its agreement with the original; a ``character specialist'' considered the choice of Chinese characters; finally a ``polisher'' made the sentences fluent. A scripture passed through several hands and several pairs of eyes before its text was finalised.

Now lean your ear into this room and listen to an argument underway. One Sanskrit word keeps appearing in the manuscript on the desk: prajñā. It roughly means a wisdom that penetrates the truth of worldly existence. How should it be translated?

Someone proposes zhihui, ``wisdom'': every Chinese reader understands it. An objection follows immediately. No: zhihui is too broad. Cleverness, quick-wittedness, and cunning can all count as wisdom. The word cannot bear the weight of prajñā; it would narrow and flatten the meaning.

What then? The eventual decision is not to translate the meaning but to copy the sound, writing bore, pronounced bō rě. A combination of syllables meaningless to Chinese people at the time is deliberately created, leaving generations of readers to fill it with meaning themselves.

Xuanzang, his predecessors, and his successors turned such choices into principles. Later writers summarised the ``five cases of non-translation'' associated with his school: five circumstances in which sound should be copied rather than meaning translated, including when a word has several meanings, when China has no corresponding thing, or when solemnity is required. In other words, after repeated arguments, the people in this hall who knew the two languages best reached a conclusion: \textbf{sometimes the best translation is not to translate.}

Imagine a young monk's bewilderment. Did we not bring these scriptures back over thousands of miles precisely so people could understand them? Why does translating them leave so many words untranslated?

Do not hurry past that puzzle. Beside the cola can in your imagination now stands a Tang oil lamp. More than a thousand years separate the two objects, but they ask the same question.

\section{Carrier or Rewriter?}
Put that question on the table without rushing to answer it.

Most people carry a default picture of translation that we might call the \textbf{transport model}: meaning is cargo; language is a box. The original meaning sits in language A's box. The translator takes it out, packs it into language B's box, and delivers it. Under this model there is one standard for good translation: nothing missing, nothing damaged. Ideally the translator is invisible, like a conveyor belt. Nobody applauds a conveyor belt, nor should anyone award it a prize.

This model sounds simple and honourable. Yet each of our two scenes has already struck it a blow.

The cola can's blow: if translation is transport, Kekou Kele is a brazen substitution. The cargo, coca leaves and kola nuts, has been discarded and replaced with the translator's own additions, good taste and happiness. Yet it is one of history's most successful translations. The translation hall's blow: if translation is transport, bore is a refusal to deliver. The cargo is too big and the box too small, so the carrier announces, ``I will not move this item; the recipient must collect it.'' Yet that is the considered choice of outstanding translators.

Another picture emerges: the \textbf{rewriting model}. Translation does not move cargo at all. After understanding the original, the translator remakes the work in another language. Think of a dish: the original supplies the ingredients, but the heat, knife work, and presentation belong to the translator. Strictly speaking, every translated novel you read in Chinese has been written again in Chinese by its translator.

This model explains much more, but its difficulties are much greater too.

If translation is rewriting, \textbf{who authorised the translator to rewrite?} Changing one word in a scripture might redirect millions of people's beliefs. Mistranslating one word in a treaty might change a border between countries. Getting one detail wrong in a detective novel might identify the wrong murderer. Calling a translator a ``carrier'' limits that person's power. Once we admit that the translator ``rewrites'', we face an uncomfortable fact: the ``foreign world'' we read is largely a version chosen for us by a handful of translators.

Distinguish several questions tangled together here. What happened belongs to evidence: Xuanzang's translation hall left the word bore, and the cola can bears Kekou Kele. Those are facts. Why it happened belongs to explanation: linguistic differences, cultural differences, and translators' considerations admit different accounts that we will compare. How people understood it at the time is another question: Xuanzang's hall regarded ``non-translation'' as caution; internet commenters today might call it laziness. How we evaluate it is another matter again. You can certainly judge bore a good or bad translation, but that is a judgement requiring reasons, not the fact itself.

Now we can state the real question: \textbf{can meaning move house across languages?} If so, what moves? If not, what exactly are translators doing every day?

\section[Literal or Free? Answers Worldwide]{Literal or Free? Answers from Around the World}
Two camps have long stood on either side of this question. Let us first give both arguments in full. Presenting only one would turn this chapter into a leaflet.

\textbf{Position one: stay close to the original; fidelity is the translator's duty.}

Let us give this camp its strongest case. It is often mocked for ``wooden translation'' and ``translationese'', which is unfair.

First, fidelity respects the author. The author wrote these words, not those. A translator who casually changes awkward wording is like an editor rewriting your composition because it reads awkwardly. The revision may indeed be better, but it is no longer yours. Second, fidelity is essential in high-stakes texts: contracts, treaties, medicine instructions, and judgements. In these, ``close enough'' means an accident. An ingredients list cannot become advertising copy. Third, and most profoundly, \textbf{staying close to the original can replenish a language.} Everyday Chinese expressions such as shijie (world), chana (instant), yinguo (cause and effect), juewu (awakening), and fangbian (convenience) were all coined or borrowed through Buddhist translation. Had Xuanzang and his colleagues pursued only smoothness, none would exist. Many expressions that seem natural in a language slowly steeped out of generations of ``awkward literal translations''. Literal translators guard not only the original but the language's future.

\textbf{Position two: meaning does not live inside words; word-for-word matching is often the least faithful translation.}

This camp's reasons are equally strong. Words never correspond one to one. English uncle forces a Chinese translator to choose immediately between a father's elder brother, his younger brother, a mother's brother, a father's sister's husband, and a mother's sister's husband. The translator must add information absent from the original: fidelity or distortion? Sentences correspond even less neatly. Japanese and Korean have elaborate honorific systems that announce the speakers' relative status as soon as they speak. Chinese lacks that machinery. A word-for-word translation either becomes unnatural speech or loses half the meaning. Meaning lives in a whole language's usage, not in individual verbal shells. Moving those shells one by one is precisely how meaning gets lost. The Kekou Kele camp would say: we discarded the shells but preserved what this product truly wants to say---it tastes good and makes you happy.

Looking around the world reveals that this was never a private concern of Chinese translation halls.

Around the ninth century, Abbasid Baghdad launched a translation movement lasting more than a hundred years. Scholars translated Greek philosophy, medicine, and mathematics into Arabic in great quantities. Payment was so generous that legend says manuscripts earned their weight in gold. Without these versions, many ancient Greek works might not have survived. Centuries later, Europeans translated them again, from Arabic into Latin. Knowledge survived through two legs of a translation relay. Transport or rewriting? It looks more like rescue.

In sixteenth-century Germany, Martin Luther translated the Bible into German. The temptation was scholarly, Latin-flavoured sentences, which would sound solemn and dependable. But he wrote a defence explaining, in essence, that translators should ask how mothers at home, children in the street, and ordinary people at market would speak: millers and weavers must understand. He chose living language and paid for it with accusations of distorting scripture.

In nineteenth-century Japan, Meiji scholars had to fit batches of Western concepts into an East Asian language. Nishi Amane coined tetsugaku, written with the characters for ``wisdom'' and ``learning'', for philosophy; science became kagaku. These coinages later travelled back to China on a large scale, as zhexue and kexue among others, and today's conversation cannot do without them. Here translation was outright invention.

Consider an extreme Chinese case. Lin Shu, in the late Qing, knew not a word of any foreign language, yet ``translated'' more than 180 foreign novels. His method was for a foreign-language collaborator to tell him the story aloud while he wrote it in elegant literary Chinese. Abridgement, rewriting, and distortion were all present. His version of \emph{The Lady of the Camellias} came from this workshop. Under the transport model it scarcely qualifies as translation. Yet these ``unqualified'' versions made a generation of Chinese readers unfamiliar with foreign literature weep over other countries' stories for the first time, profoundly influencing modern literature. This example treats a particular assumption: \textbf{the closer a translation is to its original, the greater its impact.} Lin Shu was far from the original and close to his readers.

Finally, visit a place where mistranslation can kill: the diplomatic negotiating table. In 1956, Soviet leader Nikita Khrushchev said something to Western diplomats at a reception. The interpreter rendered it in English as ``We will bury you''. The Western world erupted. Taken as a naked threat of war, the words added fuel during the Cold War's most tense years. Khrushchev later said it was a Russian idiom closer to ``We will outlive you'' or ``We will attend your funeral'': a boast that his system would have the last laugh, not an announcement of violence. Literal fidelity, disastrous effect. In 1977, during President Jimmy Carter's visit to Poland, a succession of errors by the American interpreter supplied translation classes with a lasting cautionary example. According to widespread contemporary reports, ``I left the United States this morning'' became ``I left the United States, never to return'', while remarks about hopes for the future became a version that embarrassed the whole audience.

Diplomatic interpreters occupy perhaps the world's most stressful chair. Translate literally and you may create a crisis; translate freely and you may exceed your authority by changing a leader's tone. Every second, they must carry both camps' arguments at once.

\section[Thinking Bridge: A Translation Checkup]{Thinking Bridge: Giving a Translation a Checkup}
After two thousand years of argument, can we take away a tool without choosing sides? Yes. Translation researchers' working methods suggest that instead of vaguely asking whether a version is ``accurate'', we should separate five questions. Next time you encounter a translation in a textbook, film subtitle, or software output, run it through this five-layer checklist.

\textbf{Layer one: semantics.} Has literal information been lost or altered? Turning ``I left the United States this morning'' into ``I left the United States, never to return'' fails here. This is the checklist's foundation, and the main measure for ingredients, instructions, and legal provisions.

\textbf{Layer two: rhythm and sound.} Sentences have body heat and a pulse. A succession of short, abrupt sentences conveys tension; a dragging long sentence suggests hesitation or display. A translation can get every meaning right but combine the short sentences into long ones and lose the tension. Poetry takes this further: sound is part of the content. This layer asks not ``Did you translate it correctly?'' but ``Does it still have the same heartbeat?''

\textbf{Layer three: puns and wordplay.} Translation suffers heavily here: almost every pun breaks at the seam between languages. Consider the famous lines from Tang poet Liu Yuxi's \emph{Bamboo Branch Songs}:

\begin{quote}
The sun rises in the east; rain falls in the west.\newline
They say there is no sunshine, yet sunshine remains.
\end{quote}
The charm lies in the homophones qing, ``clear weather'', and qing, ``affection''. On the surface, the weather is partly sunny and partly rainy; underneath, someone seems both to love and not to love the speaker. Translate it into any foreign language and the changed sound undoes the fastening, leaving an ordinary weather forecast. Puns cannot be moved simply by ``trying harder'': they are welded to one language's sounds.

\textbf{Layer four: cultural background.} An entire culture stands behind a word. Chinese long is auspicious; translate it as dragon and an English reader imagines a fire-breathing, treasure-hoarding monster waiting for a hero to kill it. One word, two species. In the other direction, literal translations of ``Noah's Ark'' and ``Pandora's box'' bewilder readers unfamiliar with the allusions. This layer asks whether the images called up in the translation reader's mind match those in the original reader's mind.

\textbf{Layer five: genre.} Every form has rules and expectations. Contracts require exactness; poems leave space; on-screen comments must be short; eulogies should show restraint. Translate a poem ``accurately'' into prose and you may score full marks for semantics but zero for genre, like faithfully reciting a song as a textbook lesson: all the words remain, but the song is gone. The question is whether the translated thing still counts, in its new language, as the original kind of thing.

Use the checklist by resisting immediate praise or condemnation. Ask layer by layer where a translation fails and where it succeeds. One pattern appears immediately: \textbf{no translation earns full marks on all five layers at once}. Translators always save some layers and sacrifice others. Opposing critics often talk past one another because one watches semantics while the other watches rhythm.

You can now see where translation software most easily stumbles. It is increasingly strong at semantics, but often remains oblivious to puns and culture. We will return to that shortly.

\section{Two Fates for One Verse}
Now read a short piece of primary material. It compresses all our arguments into a few dozen characters.

The \emph{Diamond Sutra} ends with a famous verse summing up the scripture's view of impermanence. Two Chinese translations survive, both by the great scripture translators mentioned earlier, both actual transmitted texts.

First, Kumarajiva's version, translated in Chang'an in the early fifth century:

\begin{quote}
All conditioned things, like dreams, illusions, bubbles, shadows, like dew and like lightning: thus should they be viewed.
\end{quote}
Now Xuanzang's version, from the seventh century, over two hundred years later:

\begin{quote}
All things made by combination, like stars, blurred vision, lamps, illusions, dew, bubbles, dreams, lightning, clouds: thus should they be viewed.
\end{quote}
They say the same thing: whatever comes together through causes and conditions is fleeting and should be seen that way. But compare them word by word: the differences are systematic.

The Sanskrit original has nine comparisons, like nine successive photographs. Xuanzang retains every one: stars, visual obscuration (an obstruction in the eye that makes nonexistent things appear), lamps, illusions, dew, bubbles, dreams, lightning, and clouds. Nine things that ``exist but cannot be relied upon''. Kumarajiva cuts them down to six: dreams, illusions, bubbles, shadows, dew, and lightning. He removes three and rearranges the Sanskrit order to make the translation smooth and rhythmic.

Which is ``right''? Apply the five-layer checklist. Semantically, Xuanzang is more complete: the subtle distinctions between the nine comparisons each carry meaning, and losing one loses a face of impermanence. Rhythmically, Kumarajiva wins decisively. His six images pair off and move as naturally as breathing, almost born for memorisation. Culturally, ``dreams, illusions, bubbles, shadows'' became a Chinese idiom, merging into the language itself, something Xuanzang's version did not achieve.

History's verdict is intriguing. For more than a thousand years, Kumarajiva's version has been the one recited throughout Chinese Buddhist temples. Xuanzang's sits in the canon, mainly consulted by scholars. The ``fuller'' version lost; the ``smoother'' version won.

Does fidelity therefore necessarily fail? No. We know today that the original had nine comparisons, and can criticise Kumarajiva for cutting three, precisely because ``clumsy'' translations such as Xuanzang's kept a record. \textbf{Fluent versions circulate; faithful versions preserve the archive. A culture needs both.}

On this dilemma, the late-Qing translator Yan Fu offered a much-overquoted but still useful formulation. In his prefatory ``General Remarks on Translation'' for \emph{Tianyan Lun}, his rendering of \emph{Evolution and Ethics}, he admitted:

\begin{quote}
Translation has three difficulties: fidelity, intelligibility, and elegance.
\end{quote}
Fidelity means honesty towards the original; intelligibility means readers can follow; elegance means the writing stands on its own. Yan Fu knew these demands often clash. For intelligibility and elegance, his own \emph{Tianyan Lun} operated extensively on the original, turning Huxley's scientific prose into Tongcheng-school literary Chinese. Placing the three aims together is an honest confession: translation offers no all-winning option, only choices about what you will sacrifice.

\section[Why Do Translations Differ?]{Why Do Translations of One Original Differ So Much?}
Tighten the question another turn. Versions of the same passage can sometimes seem like different works. Why? Here are three genuinely different explanations, each reasonable and each limited. Remember that we are explaining the origin of differences: this ``why'' concerns explanation, not a dispute over facts.

\textbf{Explanation one: languages have different structures.}

This is the most intuitive answer: differences are forced by different linguistic machinery. Vocabulary boundaries fail to align, as uncle splits into five Chinese relatives. Grammatical mechanisms fail to align: honorifics, tense, singular and plural. Sound inventories fail to align, killing puns. On this account the translator dismantles machinery; differences are technical losses, and the job is to minimise them. Clear and practical, this explains most individual examples. Its limit is that it cannot explain why versions between \textbf{the same pair of languages} differ so much across eras. The \emph{Diamond Sutra}'s Sanskrit did not change, nor did the broad framework of Chinese, yet Kumarajiva and Xuanzang produced different works. If the machine stayed the same but the product changed, machinery cannot explain everything.

\textbf{Explanation two: each age has its own idea of a ``good translation''.}

This shifts attention from language to history. When Wei and Jin readers first encountered Buddhist texts, a method called geyi, ``matching meanings'', became popular: concepts from Laozi and Zhuangzi were compared with Buddhist ones, using benwu, ``original nonbeing'', to translate Buddhist ``emptiness''. These were the tools available in Chinese readers' intellectual toolbox, so they were borrowed first. Yan Fu translated science into Tongcheng literary prose because his intended readers were scholar-officials trained in eight-legged essays; plain-language science carried little weight with them. Today, the same book would probably become clear vernacular prose, perhaps with an internet flavour. On this account, a translation is shaped by its era's aesthetics and readers' tastes; the translator executes the age's preferences. It explains what the first account cannot: an unchanged machine yielding a changed product. Its limit is making the translator an era's puppet. Kumarajiva and Xuanzang faced broadly the same China but made opposite choices, and every age has translators who resist the current. An era does not explain individual differences.

\textbf{Explanation three: translation has a task, and the task determines the method.}

This asks different questions: who commissioned it, who will read it, and why? Luther wanted ordinary believers to understand scripture directly, so he naturally chose living speech. The English Bible authorised by King James I was meant to unify church readings nationwide, so it sought solemnity, authority, and memorable sound. A drinks company wants sales, so Kekou Kele places advertising effect above lexical meaning. Today's fan-subtitling groups serve fellow enthusiasts and freely use internet slang. Here differences arise from commissions and intended audiences. Before asking ``Is it well translated?'', ask ``What is this translation for?'' This sharp account sees motives the others miss. Its limit is that attributing everything to tasks and interests can overlook what translators seek without calculation: reverence for the original, obsession with a word, or a craftsperson's simple desire to get a sentence right. Not every choice reduces to a purpose.

None of the three accounts can swallow every difference. That is not embarrassing but instructive: \textbf{when an explanation claims to explain everything, it has probably thrown something away in advance.} For an actual passage, the most honest approach uses all three: how much was forced by language, shaped by the age, or determined by the task?

\section[Further Challenge: Translation as Rewriting]{Further Challenge: Translation as Rewriting}
Now bring in the chapter's weightiest theory. In the later twentieth century, scholars established ``translation studies'' as a discipline. One, Belgian-born André Lefevere, wrote a book whose title already answers our question: \emph{Translation, Rewriting, and the Manipulation of Literary Fame}. His core claim fits one sentence: \textbf{translation is a form of rewriting.}

Lefevere argued that translators never merely move words one to one. Three nets hang over them. The first is poetics: what their age calls ``good literature'' and ``good translation'', expectations difficult to resist. Kumarajiva's fluency and Yan Fu's classical prose are products of their ages' poetics. The second is patronage: the powers that fund, authorise, publish, and distribute translations, including courts, churches, publishers, drinks companies, and platforms. They decide what is translated, how, and for whom. The third is ideology: what should be strengthened, softened, or removed. Translators move among these nets, and every word bears their impressions. The ``foreign literature'' you read is never the foreign world itself but a rewriting passed through three filters.

Does that make translators too passive? In a famous early-nineteenth-century lecture, the German scholar Schleiermacher described their choices more actively. In essence, he said there are only two routes: \textbf{leave the author as undisturbed as possible and lead the reader towards the author, or leave the reader as undisturbed as possible and lead the author towards the reader.} In the late twentieth century, American scholar Lawrence Venuti gave these routes their now-common names. Bringing the author over is \textbf{domestication}: the translation reads so smoothly that it seems not to be translated, as if the author could write Chinese. Sending the reader over is \textbf{foreignisation}: deliberately retaining the original's unfamiliarity, making the reader stumble slightly and recognise something foreign.

Kekou Kele is extreme domestication: zero unfamiliarity, a fully localised product. Xuanzang's bore and pusa, ``bodhisattva'', are extreme foreignisation: better invent transliterations nobody understands than lose the foreign concepts' weight. Kumarajiva's verse domesticates; Xuanzang's foreignises. These concepts turn two millennia of muddled argument over literal versus free translation into identifiable choices: this time, whom did the translator move towards whom?

Venuti makes a sharper point. The smoother a domesticated translation reads, the less visible its translator becomes. Readers think they are reading the author directly; the translator's labour, decisions, and power vanish behind fluent sentences. He calls this ``the translator's invisibility''. Once the translator disappears, so does one culture's remaking of another. British and American readers accustomed to smooth domesticated versions can mistakenly suppose all world literature naturally suits their tastes. Unequal cultural relations stand behind that fluency. Now you see why we began with transport or rewriting. For Lefevere and Venuti, translation always rewrites; the question is whether \textbf{the traces of rewriting are concealed or revealed}.

But grasp the theory's other half or you may fall into a new extreme. Learning about ``untranslatability'' and ``rewriting'' can produce fashionable despair: nothing can be translated accurately, puns inevitably die, cultural barriers are mountains, so speakers of different languages can never truly understand one another. That is a new superstition born of excessive theorising. \textbf{Untranslatable does not mean incomprehensible.} Distinguish ``cannot be said again in the same way'' from ``meaning is locked away forever''. They are entirely different. The pun on sunshine and affection cannot be transported intact, but one explanation---``the Chinese words for clear weather and affection sound alike''---lets readers of any language understand its mechanism within seconds. They lose the pleasure of the sound's immediate springing into action but gain a full understanding of how it works. Untranslatable sound can be explained; an untranslatable rhythm's heartbeat can be described. Human understanding across languages has never depended simply on moving sounds about. It also relies on explanations, annotations, examples, and analogies: all the patient communicative work beyond translation. Those of us raised on translations can still be moved by haiku and weep over \emph{Les Misérables}. That is evidence.

The theory draws a bright line: translation rewrites, and rewriting is constrained by nets. Yet beyond it remains something theory cannot collect away: people's curiosity about other people is older than any net.

\section{Babel on Your Phone}
This two-thousand-year-old problem is happening millions of times a second inside your pocket.

Machine translation, entrusting translation to computers, has passed through several generations. Early systems calculated rigidly from human-written grammar rules and dictionaries, often producing unnatural language. Statistical methods followed, counting probabilities in vast collections of existing translations. Today's mainstream neural machine translation and large language models learn linguistic correspondences from enormous amounts of text. Their effectiveness is real. Abroad, a photograph can help you read a menu, ask directions, or understand product instructions. Running foreign materials through a machine before tackling the original can double efficiency. Multinational customer service and in-game chat depend on it. Machines have taken over much everyday semantic work.

Where do they fail? The five-layer checklist makes the trouble clear.

Irony and sarcasm: a machine misses the discrepancy and faithfully translates a sarcastic ``You're really something'' as praise. Puns: as we saw, it may not even realise a pun exists, choosing one meaning without expression. Honorifics and status: it often misplaces the machinery of Japanese and Korean social relations. Cultural gaps: it does not know that a monster stands between Chinese long and English dragon.

High-stakes texts are the most serious case. Medical and legal translation errors are documented. In a famous American medical incident in the 1980s, a Spanish-speaking patient's family described him as intoxicado. In Spanish the word can broadly suggest having eaten something harmful and feeling sick, but people at the scene understood it through English intoxicated, meaning drunk or drug-overdosed. Treatment went in the wrong direction, the patient suffered lifelong paralysis, and the hospital ultimately paid tens of millions of dollars. Two faces of one word separated two diagnoses. Contracts, medical records, and judgements are still not confidently entrusted to machines alone.

One easily misjudged counterexample deserves emphasis. You might think machine translation's greatest danger is bad writing: broken sentences obviously look false. Wrong. Bad writing is comparatively safe because it visibly says ``Do not trust me'', encouraging verification. \textbf{The real danger is fluent error.} Today's machine mistakes often come in smooth, confident, professional sentences, errors dressed as correctness. In the age of obvious translationese, errors exposed themselves; in the AI age, they hide. The first rule for reading machine output is therefore: the smoother it is, the more alert you must be to cargo quietly exchanged where you cannot see. Remember, a machine is a ``rewriter'' too, and never announces when it rewrites.

If machines are so capable, is learning a foreign language still useful? Do not rush to answer. A machine can translate the literal words of a friend's letter, but ``Why choose this particular word?'', ``Is this politeness or anger?'', and ``Should I respond to this joke?'' belong to the checklist's upper layers. Machines offer hints, not certainty. Learning a language has never meant only vocabulary and grammar. It provides entry to the culture behind it: how its puns make people laugh, what its silences mean, what its ``dragon'' looks like. A machine can deliver the cargo, but inspecting it still requires knowledge of the language.

\section[Whom Would You Trust with a Letter?]{For You: Whom Would You Trust with a Letter?}
End with a concrete dilemma.

Suppose someone especially important to you is about to spend a long time in a country whose language they do not know. Before departure, you want to write a letter, perhaps a belated apology or something never said aloud. For it to be understood, the letter must become another language. You have three options:

A. Translate it yourself, word by word, with a dictionary. This stays closest to your intentions but may be stiff or even ridiculous.

B. Give it to translation software. This is smoothest and easiest, but it quietly rewrites where you cannot see and never announces it.

C. Ask a fluent bilingual friend. The translation is good and accurate, but first the letter passes through a third person's heart. They choose your tone and decide which sentences deserve emphasis and which should be lighter.

Which do you choose? Give reasons and persuade someone who chose differently.

Before deciding, weigh what you now hold. You know why word-for-word matching often fails: words do not align, puns are welded to sound, and cultures attach different pictures to the same word. You have seen Xuanzang's translation hall, the two versions of the verse, Luther's choice, and the turmoil around Khrushchev's words. You have the five-layer checklist, domestication and foreignisation, the theory of translation as rewriting, and its antidote: untranslatable does not mean incomprehensible. You also know the greatest danger is not poor translation but fluent error.

None of that answers the innermost question for you. What must this letter not lose? Its literal words, its tone, the instant when its pun clicks, or the fact that ``these words really came from my hand''? How much rewriting will you accept to preserve it?

Is translation transport or rewriting? Perhaps your answer is: it depends on the cargo and whose hands you trust. Then whose hands do you trust, and why?
